% !TEX program = xelatex
\documentclass[12pt,a4paper]{ctexart}

\usepackage{geometry}
\geometry{margin=2.5cm}

\usepackage{listings}
\usepackage{xcolor}
\usepackage{hyperref}
\usepackage{booktabs}
\usepackage{enumitem}
\usepackage{fancyvrb}
\usepackage{tcolorbox}
\tcbuselibrary{listings,skins,breakable}

\hypersetup{
  colorlinks=true,
  linkcolor=blue!70!black,
  urlcolor=blue!70!black,
}

% 代码样式
\lstset{
  basicstyle=\ttfamily\small,
  backgroundcolor=\color{gray!8},
  frame=single,
  rulecolor=\color{gray!30},
  breaklines=true,
  columns=fullflexible,
  keepspaces=true,
  showstringspaces=false,
  commentstyle=\color{green!50!black},
  keywordstyle=\color{blue!70!black}\bfseries,
  stringstyle=\color{red!60!black},
}

\lstdefinelanguage{slurm}{
  morekeywords={SBATCH,sbatch,srun,squeue,scancel,seff,sacct,scontrol,module,load,conda,activate,python},
  morecomment=[l]{\#},
  morestring=[b]",
}

\newtcolorbox{tipbox}[1][]{
  colback=blue!5,colframe=blue!50!black,
  fonttitle=\bfseries,title=提示,#1
}

\newtcolorbox{warnbox}[1][]{
  colback=red!5,colframe=red!50!black,
  fonttitle=\bfseries,title=注意,#1
}

\newtcolorbox{conceptbox}[1][]{
  colback=green!5,colframe=green!50!black,
  fonttitle=\bfseries,title=核心概念,#1
}

\title{\textbf{Tufts HPC 集群使用教程}\\[6pt]
\large 从零开始，以强化学习实验矩阵为例}
\author{}
\date{}

\begin{document}
\maketitle
\tableofcontents
\newpage

% ---- 第一部分：速查与快速上手（每天都用）----
%=============================================================================
\section*{日常速查表}
\addcontentsline{toc}{section}{日常速查表}
%=============================================================================

以下是你在集群上最常用的操作，不需要翻教程正文，直接复制粘贴即可。

\subsection*{连接与节点}

\begin{table}[h]
\centering
\small
\begin{tabular}{lp{9cm}}
\toprule
操作 & 命令 \\
\midrule
SSH 登录 & \texttt{ssh your\_utln@login-prod.pax.tufts.edu} \\
快速开一个 CPU 计算节点 & \texttt{srun -p batch -t 0-2:00:00 -n 2 --mem=4g --pty bash} \\
快速开一个 GPU 计算节点 & \texttt{srun -p gpu -t 0-2:00:00 -n 2 --mem=4g --gres=gpu:1 --pty bash} \\
退出计算节点 & \texttt{exit} \\
\bottomrule
\end{tabular}
\end{table}

\begin{tipbox}
很多操作（查目录大小、安装 Conda 环境等）在登录节点上会很慢，\textbf{开一个 CPU 计算节点再操作会快很多}。
\end{tipbox}

\subsection*{作业管理}

\begin{table}[h]
\centering
\small
\begin{tabular}{lp{9cm}}
\toprule
操作 & 命令 \\
\midrule
提交作业 & \texttt{sbatch script.sh} \\
查看我的作业 & \texttt{squeue --me} \\
取消某个作业 & \texttt{scancel JOBID} \\
取消我的所有作业 & \texttt{scancel -u \$USER} \\
查看已完成作业的效率 & \texttt{seff JOBID} \\
\bottomrule
\end{tabular}
\end{table}

\subsection*{资源与存储}

\begin{table}[h]
\centering
\small
\begin{tabular}{lp{9cm}}
\toprule
操作 & 命令 \\
\midrule
查看我的配额（实时刷新） & \texttt{watch -n 5 quota -s} \\
查看 Home 目录各文件夹大小 & \texttt{du -sh \textasciitilde/* \textasciitilde/.* 2>/dev/null | sort -rh | head -20} \\
查看集群 GPU 空闲情况 & \texttt{module load hpctools \&\& hpctools}（选 2） \\
查看各分区状态 & \texttt{sinfo} \\
查看 GPU 使用（在计算节点上） & \texttt{nvidia-smi} \\
\bottomrule
\end{tabular}
\end{table}

\subsection*{文件传输（在本地电脑上运行）}

\begin{table}[h]
\centering
\small
\begin{tabular}{lp{9cm}}
\toprule
操作 & 命令 \\
\midrule
上传项目目录 & \texttt{rsync -azP ./project/ user@login-prod.pax.tufts.edu:\textasciitilde/project/} \\
下载结果 & \texttt{rsync -azP user@login-prod.pax.tufts.edu:\textasciitilde/project/results/ ./results/} \\
\bottomrule
\end{tabular}
\end{table}

\subsection*{环境管理}

\begin{table}[h]
\centering
\small
\begin{tabular}{lp{9cm}}
\toprule
操作 & 命令 \\
\midrule
加载 Conda & \texttt{module load miniforge/25.3.0} \\
激活环境 & \texttt{conda activate rl\_env} \\
清理 uv 缓存（容易撑满配额） & \texttt{rm -rf \textasciitilde/.cache/uv} \\
清理 pip 缓存 & \texttt{rm -rf \textasciitilde/.cache/pip} \\
清理 conda 包缓存 & \texttt{conda clean --all -y} \\
\bottomrule
\end{tabular}
\end{table}

\newpage
          % 日常速查表
%=============================================================================
\section{Quick Start：科研工作者的日常流程}
%=============================================================================

如果你只想尽快跑起来，这一节就够了。后面的章节是详细解释。

\subsection{一次典型的工作流}

假设你已经在本地写好了代码，想在集群上跑。整个过程只需要 5 步：

\textbf{第 1 步：把代码同步到集群}

在你的\textbf{本地电脑}上运行：

\begin{lstlisting}[language=bash]
rsync -azP --exclude='.venv' --exclude='__pycache__' \
  ./my_project/ your_utln@login-prod.pax.tufts.edu:~/my_project/
\end{lstlisting}

如果你已经配置了 SSH 别名（见连接章节），可以简写为：

\begin{lstlisting}[language=bash]
rsync -azP --exclude='.venv' --exclude='__pycache__' \
  ./my_project/ hpc:~/my_project/
\end{lstlisting}

\texttt{rsync} 只传有变化的文件，所以第二次开始会很快。

\textbf{第 2 步：SSH 登录集群}

\begin{lstlisting}[language=bash]
ssh hpc    # 或 ssh your_utln@login-prod.pax.tufts.edu
\end{lstlisting}

\textbf{第 3 步：提交作业}

假设你的项目里已经有一个 SLURM 脚本 \texttt{run.sh}：

\begin{lstlisting}[language=bash]
cd ~/my_project
sbatch run.sh
# 输出: Submitted batch job 347502
\end{lstlisting}

不知道怎么写 SLURM 脚本？这里是一个最小的模板，保存为 \texttt{run.sh}：

\begin{lstlisting}[language=slurm]
#!/bin/bash
#SBATCH -J my_job               # 作业名
#SBATCH --time=00-06:00:00      # 最多跑 6 小时
#SBATCH -p gpu                  # 使用 GPU 分区
#SBATCH --gres=gpu:1            # 1 块 GPU
#SBATCH --cpus-per-task=4       # 4 个 CPU 核
#SBATCH --mem=16g               # 16GB 内存
#SBATCH --output=%x.%j.out      # 标准输出日志
#SBATCH --error=%x.%j.err       # 错误日志

module load miniforge/25.3.0    # 加载 Conda
module load cuda/12.9.0         # 加载 CUDA
conda activate my_env           # 激活你的环境

python train.py                 # 运行你的程序
\end{lstlisting}

如果不需要 GPU，把 \texttt{-p gpu} 改成 \texttt{-p batch}，删掉 \texttt{--gres} 那行。

\textbf{第 4 步：查看作业状态}

\begin{lstlisting}[language=bash]
squeue --me                    # 看作业是在跑(R)还是排队(PD)
cat my_job.347502.out          # 看输出日志（347502 换成你的 Job ID）
\end{lstlisting}

提交之后你可以断开 SSH 去做别的事，作业会在后台继续跑。

\textbf{第 5 步：下载结果}

作业跑完后，在你的\textbf{本地电脑}上运行：

\begin{lstlisting}[language=bash]
rsync -azP hpc:~/my_project/results/ ./results/
\end{lstlisting}

\subsection{每天重复的部分}

上面 5 步里，第 2-4 步是你每天都会做的。最常见的循环是：

\begin{enumerate}[nosep]
  \item 本地改代码
  \item \texttt{rsync} 同步到集群（10 秒）
  \item \texttt{ssh hpc}，\texttt{sbatch run.sh}（5 秒）
  \item 过一会儿 \texttt{squeue --me} 看状态
  \item 跑完后 \texttt{rsync} 拉结果回来
\end{enumerate}

\subsection{遇到问题时怎么办}

\begin{table}[h]
\centering
\small
\begin{tabular}{lp{8cm}}
\toprule
情况 & 怎么做 \\
\midrule
作业一直排队（PD） & \texttt{sinfo} 看分区是否有空闲节点；减少资源请求；尝试加 \texttt{preempt} 分区 \\
作业运行失败 & 看 \texttt{.err} 日志文件找报错信息 \\
不确定环境对不对 & 开一个交互式节点调试：\texttt{srun -p batch -t 0-1:00:00 --mem=4g --pty bash} \\
配额满了 & \texttt{watch -n 5 quota -s} 监控；\texttt{du -sh \textasciitilde/.cache/*} 找大文件；\texttt{rm -rf \textasciitilde/.cache/uv} 清缓存 \\
想取消所有作业 & \texttt{scancel -u \$USER} \\
\bottomrule
\end{tabular}
\end{table}
         % Quick Start：5 步跑起来

% ---- 第二部分：常用操作详解（经常需要翻阅）----
%=============================================================================
\section{连接到集群}
%=============================================================================

\subsection{新旧集群：确保你连的是新集群}

Tufts 目前有两套 HPC 集群——\textbf{旧集群}正在逐步淘汰，\textbf{新集群}（升级版）是官方推荐使用的。两者的登录地址不同：

\begin{table}[h]
\centering
\begin{tabular}{lll}
\toprule
 & SSH 登录地址 & OnDemand 网址 \\
\midrule
\textbf{新集群}（推荐） & \texttt{login-prod.pax.tufts.edu} & \texttt{ondemand-prod.pax.tufts.edu} \\
旧集群（淘汰中） & \texttt{login.pax.tufts.edu} & \texttt{ondemand.pax.tufts.edu} \\
\bottomrule
\end{tabular}
\end{table}

\begin{warnbox}
\textbf{请始终使用新集群}（带 \texttt{-prod} 后缀的地址）。旧集群仍然可以登录，但 Tufts 官方强烈建议迁移到新集群，旧集群将来会被关闭。如果你之前的脚本或配置里用了 \texttt{login.pax.tufts.edu}（不带 \texttt{-prod}），请尽快更新。
\end{warnbox}

本教程所有内容均基于\textbf{新集群}。

\subsection{SSH 登录}

通过 SSH 连接到新集群的登录节点：

\begin{lstlisting}[language=bash]
ssh your_utln@login-prod.pax.tufts.edu
\end{lstlisting}

这里 \texttt{your\_utln} 是你的 Tufts 用户名（全小写）。如果你在校外，需要先连接 Tufts VPN（或使用跳板机，见后文）。

登录成功后，你会看到类似这样的提示符：

\begin{lstlisting}
[your_utln@login-p01 ~]$
\end{lstlisting}

\texttt{login-p01} 表示你当前在第 1 个登录节点上。这个编号不重要，3 个登录节点看到的文件完全一样。

\subsection{OnDemand 网页界面}

如果你不熟悉命令行，也可以通过浏览器访问：

\begin{center}
\url{https://ondemand-prod.pax.tufts.edu/}
\end{center}

OnDemand 提供图形化的文件管理器、终端、以及 Jupyter Notebook 等应用，适合入门和调试。

\subsection{进阶：免密码登录与跳板机（无需 VPN）}

每次登录都要输密码很烦。你可以配置 \textbf{SSH Key} 实现免密登录。

\subsubsection{第一步：生成密钥对（如果还没有）}

在你的\textbf{本地电脑}上运行：

\begin{lstlisting}[language=bash]
# 检查是否已有密钥
ls ~/.ssh/id_*.pub

# 如果没有，生成一对
ssh-keygen -t ed25519
# 一路回车即可（不设密码短语）
\end{lstlisting}

\subsubsection{第二步：把公钥传到集群上}

\begin{lstlisting}[language=bash]
ssh-copy-id your_utln@login-prod.pax.tufts.edu
\end{lstlisting}

输入一次密码后，以后就不需要了。验证：

\begin{lstlisting}[language=bash]
ssh your_utln@login-prod.pax.tufts.edu
# 应该直接进去，不再问密码
\end{lstlisting}

\begin{tipbox}
如果免密登录不生效，检查集群上的文件权限：
\begin{lstlisting}[language=bash]
chmod 700 ~/.ssh
chmod 600 ~/.ssh/authorized_keys
\end{lstlisting}
\end{tipbox}

\subsubsection{第三步（可选）：配置别名，简化命令}

每次输 \texttt{ssh your\_utln@login-prod.pax.tufts.edu} 太长了。在本地编辑 \texttt{\textasciitilde/.ssh/config}，添加：

\begin{lstlisting}
Host hpc
    HostName login-prod.pax.tufts.edu
    User your_utln
\end{lstlisting}

之后只需要 \texttt{ssh hpc} 即可登录，\texttt{rsync}、\texttt{scp} 也可以用这个别名：

\begin{lstlisting}[language=bash]
rsync -azP ./project/ hpc:~/project/
scp results.csv hpc:~/
\end{lstlisting}

\subsubsection{第四步（可选）：通过跳板机连接，不需要 VPN}

如果你在校外又不想用 VPN，可以通过一台\textbf{校内能访问的服务器}作为跳板（ProxyJump）。前提是你有一台可以从外网访问的校内机器（比如通过 ngrok 暴露的实验室服务器）。

在本地 \texttt{\textasciitilde/.ssh/config} 中配置：

\begin{lstlisting}
# 跳板机（你的校内服务器，通过 ngrok 等方式从外网可达）
Host my_jump_server
    HostName your_server_address
    User your_username
    Port your_port

# HPC，通过跳板机连接
Host hpc
    HostName login-prod.pax.tufts.edu
    User your_utln
    ProxyJump my_jump_server
\end{lstlisting}

\texttt{ProxyJump} 告诉 SSH：先连跳板机，再从跳板机连 HPC。对你来说就是一条命令：

\begin{lstlisting}[language=bash]
ssh hpc    # 自动经过跳板机，直达 HPC 登录节点
\end{lstlisting}

配合免密登录（对跳板机和 HPC 都做 \texttt{ssh-copy-id}），整个过程零密码、无 VPN。

\begin{lstlisting}[language=bash]
# 分别对跳板机和 HPC 配置免密（各输一次密码，之后永久免密）
ssh-copy-id my_jump_server
ssh-copy-id hpc
\end{lstlisting}

\begin{warnbox}
跳板机方案的前提是你有一台校内服务器可以从外网访问。如果你没有这样的机器，还是需要使用 Tufts VPN。
\end{warnbox}
             % 连接：SSH、免密、跳板机
%=============================================================================
\section{SLURM 详解：提交和管理作业}
%=============================================================================

现在你已经了解了集群的基本结构、文件存储、以及怎么准备软件环境。接下来我们来看怎么真正提交一个计算任务。

\subsection{作业脚本的结构}

一个 SLURM 作业脚本就是一个普通的 Bash 脚本，但开头多了一些 \texttt{\#SBATCH} 开头的特殊注释。SLURM 读这些注释来了解你需要什么资源。

我们来逐行看一个完整的例子：

\begin{lstlisting}[language=slurm]
#!/bin/bash
\end{lstlisting}

这行声明这是一个 Bash 脚本，是所有 shell 脚本的标准开头。

\begin{lstlisting}[language=slurm]
#SBATCH -J my_experiment       # 作业名称（方便你在队列里识别）
#SBATCH --time=00-06:00:00     # 最多运行 6 小时（格式：DD-HH:MM:SS）
\end{lstlisting}

\texttt{-J} 设置作业名称，会显示在 \texttt{squeue} 输出里。\texttt{--time} 设置时间上限——超过这个时间作业会被强制终止。\textbf{所有公共分区的时间上限是 2 天}（\texttt{2-00:00:00}）。

\begin{lstlisting}[language=slurm]
#SBATCH -p gpu                 # 使用 gpu 分区
#SBATCH --gres=gpu:1           # 需要 1 块 GPU
#SBATCH --cpus-per-task=4      # 需要 4 个 CPU 核
#SBATCH --mem=16g              # 需要 16GB 内存
\end{lstlisting}

这几行声明你需要的\textbf{硬件资源}。\texttt{-p} 选择分区（后面会详细讲），\texttt{--gres} 请求 GPU，\texttt{--cpus-per-task} 和 \texttt{--mem} 分别请求 CPU 和内存。

\begin{lstlisting}[language=slurm]
#SBATCH --output=my_exp.%j.out # 标准输出写入此文件
#SBATCH --error=my_exp.%j.err  # 错误输出写入此文件
\end{lstlisting}

这两行\textbf{非常重要}——它们决定了你的程序输出（\texttt{print} 的内容、报错信息）保存在哪里。\texttt{\%j} 会被替换为作业 ID，这样每次运行都有独立的日志文件。

\begin{lstlisting}[language=slurm]
# ---- 以下是普通的 shell 命令 ----
module load cuda/12.9.0        # 加载 CUDA
module load miniforge/25.3.0    # 加载 Conda
conda activate rl_env           # 激活环境

python train.py --lr 0.001      # 运行你的程序
\end{lstlisting}

\texttt{\#SBATCH} 行之后，就是普通的 shell 命令，和你在终端里敲的一模一样。

\begin{warnbox}
\texttt{\#SBATCH} 行看起来像注释，但 SLURM 会解析它们。它们\textbf{必须紧跟在 \texttt{\#!/bin/bash} 之后}，在任何非注释命令之前。如果中间插了一行普通命令（哪怕是 \texttt{echo}），后面的 \texttt{\#SBATCH} 都会被忽略。
\end{warnbox}

\subsection{输出文件在哪里？}

这是很多新手最常问的问题。答案取决于你怎么写 \texttt{--output} 和 \texttt{--error}：

\subsubsection{SLURM 日志文件（标准输出/错误）}

\begin{itemize}
  \item 如果你写了 \texttt{--output=my\_exp.\%j.out}，日志文件就在\textbf{你提交作业时所在的目录}下
  \item 如果你写了 \texttt{--output=logs/\%x.\%j.out}，就在那个目录的 \texttt{logs/} 子目录下（目录必须提前创建）
  \item 如果你\textbf{没写} \texttt{--output}，默认输出到 \texttt{slurm-JOBID.out}，也在提交目录下
\end{itemize}

常用的文件名变量：

\begin{table}[h]
\centering
\begin{tabular}{ll}
\toprule
变量 & 含义 \\
\midrule
\texttt{\%j} & Job ID（如 \texttt{296794}） \\
\texttt{\%x} & 作业名称（\texttt{-J} 设置的名字） \\
\texttt{\%A} & Array Job 的主 ID \\
\texttt{\%a} & Array Task 的序号 \\
\texttt{\%N} & 运行节点名 \\
\bottomrule
\end{tabular}
\end{table}

\subsubsection{你的程序自己写出的文件}

除了 SLURM 日志，你的程序通常还会写出其他文件（模型权重、CSV 结果等）。这些文件\textbf{保存在你程序里指定的路径下}。

例如，如果你的 \texttt{train.py} 里写了：

\begin{lstlisting}[language=python]
torch.save(model, "results/model.pt")
\end{lstlisting}

那么 \texttt{results/model.pt} 就在\textbf{计算节点执行脚本时的工作目录下}——通常就是你提交 \texttt{sbatch} 时所在的目录。

\begin{tipbox}
\textbf{最佳实践}：在程序里用绝对路径或相对于项目根目录的路径来保存结果，避免文件散落到意想不到的地方。例如：

\begin{lstlisting}[language=python]
output_dir = f"/cluster/tufts/mylab/myutln/project/results/{env}/{arch}/seed_{seed}"
os.makedirs(output_dir, exist_ok=True)
\end{lstlisting}
\end{tipbox}

\subsection{怎样取回结果文件？}

作业跑完后，结果文件已经在集群的共享存储上了。你有几种方式获取它们：

\textbf{方式一：直接在集群上查看}

\begin{lstlisting}[language=bash]
# SSH 登录后直接查看
cat my_exp.296794.out              # 看日志
ls results/                        # 看结果文件
\end{lstlisting}

\textbf{方式二：下载到本地（小文件用 scp）}

\begin{lstlisting}[language=bash]
# 在你的本地电脑上运行
scp your_utln@login-prod.pax.tufts.edu:~/project/results/summary.csv ./
\end{lstlisting}

\textbf{方式三：下载到本地（大量文件用 rsync）}

\begin{lstlisting}[language=bash]
# 只同步有变化的结果文件，增量下载
rsync -azP your_utln@login-prod.pax.tufts.edu:~/project/results/ ./results/
\end{lstlisting}

\textbf{方式四：超大数据用 Globus}

如果结果文件有几十 GB（比如大量 checkpoint），用 Globus 后台传输最稳定。参见第 \ref{sec:file-transfer} 节。

\subsection{提交和管理作业}

\textbf{提交作业：}

\begin{lstlisting}[language=bash]
$ sbatch my_job.sh
Submitted batch job 296794      # SLURM 返回 Job ID，记住这个号
\end{lstlisting}

\textbf{查看作业状态：}

\begin{lstlisting}[language=bash]
$ squeue --me
  JOBID PARTITION  NAME     USER  ST  TIME  NODES NODELIST
 296794       gpu  my_exp  yzhang  R  3:42      1 pax003
\end{lstlisting}

状态码含义：\texttt{R} = 正在运行，\texttt{PD} = 在排队等资源，\texttt{CG} = 即将完成。

常见排队原因：\texttt{Resources}（资源不足）、\texttt{Priority}（优先级低）、\texttt{QOSMax}（超出配额限制）。

\textbf{取消作业：}

\begin{lstlisting}[language=bash]
scancel 296794          # 取消指定作业
scancel -u $USER        # 取消你的所有作业
\end{lstlisting}

\textbf{查看已完成作业的资源效率：}

\begin{lstlisting}[language=bash]
seff 296794
\end{lstlisting}

这会显示你实际用了多少 CPU 和内存。如果利用率很低，下次可以减少请求量——\textbf{请求越精确，排队时间越短}。

\subsection{交互式会话：调试用}

如果你想在计算节点上直接敲命令（比如调试程序、测试环境），可以用 \texttt{srun} 申请一个交互式会话：

\begin{lstlisting}[language=bash]
# CPU 交互式会话
srun -p batch -t 0-2:00:00 -n 2 --mem=4g --pty bash

# GPU 交互式会话
srun -p gpu -t 0-2:00:00 -n 2 --mem=4g --gres=gpu:1 --pty bash
\end{lstlisting}

\texttt{--pty bash} 的意思是``给我一个交互式的 bash 终端''。进入后你会看到提示符从 \texttt{login-p01} 变成类似 \texttt{pax006} 的计算节点名，说明你已经在计算节点上了。用 \texttt{exit} 退出并释放资源。

\begin{tipbox}
退出时你可能会看到类似这样的报错：
\begin{lstlisting}
srun: error: pax017: task 0: Exited with exit code 1
\end{lstlisting}
\textbf{这通常是无害的}，不代表出了什么问题。它只是说你的 shell 退出时状态码不是 0——最常见的原因是 Conda 的 \texttt{(base)} 环境激活脚本或之前某个命令（比如打错了一个命令名）在 shell 内部留下了非零退出码。资源已经正常释放了，不用担心。
\end{tipbox}

\begin{tipbox}
交互式会话的 QOS 允许最多 4 小时、1 个 GPU。适合调试，不适合跑完整实验。
\end{tipbox}
       % SLURM：写脚本、提交、输出文件、交互调试
%=============================================================================
\section{文件传输：怎样上传和下载文件？}
\label{sec:file-transfer}
%=============================================================================

这是使用 HPC 的第一步，也是很多教程忽略的重要环节。

\subsection{场景一：传几个文件（用 scp）}

\texttt{scp}（Secure Copy）是最简单的文件传输命令。它的语法和 \texttt{cp} 很像，只不过可以跨机器复制。

\textbf{从本地上传到集群：}

\begin{lstlisting}[language=bash]
# 上传单个文件
scp train.py your_utln@login-prod.pax.tufts.edu:~/project/

# 上传整个目录（加 -r 参数）
scp -r my_project/ your_utln@login-prod.pax.tufts.edu:~/project/
\end{lstlisting}

冒号后面的部分是集群上的目标路径。\texttt{\textasciitilde/} 表示你的 Home 目录。

\textbf{从集群下载到本地：}

\begin{lstlisting}[language=bash]
# 下载单个文件
scp your_utln@login-prod.pax.tufts.edu:~/project/results.csv ./

# 下载整个结果目录
scp -r your_utln@login-prod.pax.tufts.edu:~/project/results/ ./local_results/
\end{lstlisting}

\begin{tipbox}
\texttt{scp} 的格式是 \texttt{scp 源路径 目标路径}，谁在前面谁就是``从哪里复制''。
本地路径直接写，远程路径格式为 \texttt{用户名@主机:路径}。
\end{tipbox}

\subsection{场景二：传一个包含几十个文件的项目（用 rsync）}

当你的项目有很多文件时，\texttt{rsync} 比 \texttt{scp} 更好用，因为它有两个重要优势：
\begin{itemize}[nosep]
  \item \textbf{增量同步}：只传输有变化的文件，不会重复传输没改过的文件
  \item \textbf{断点续传}：如果中途断了，重新运行会从断点继续，不用从头来
\end{itemize}

\textbf{上传整个项目：}

\begin{lstlisting}[language=bash]
# -a: 保留文件属性   -z: 压缩传输   -P: 显示进度
rsync -azP ./my_experiment/ \
  your_utln@login-prod.pax.tufts.edu:~/project/my_experiment/
\end{lstlisting}

\textbf{下载实验结果：}

\begin{lstlisting}[language=bash]
rsync -azP your_utln@login-prod.pax.tufts.edu:~/project/results/ ./results/
\end{lstlisting}

\textbf{排除不需要的文件}（比如本地的虚拟环境或缓存）：

\begin{lstlisting}[language=bash]
rsync -azP --exclude='.venv' --exclude='__pycache__' --exclude='.git' \
  ./my_experiment/ \
  your_utln@login-prod.pax.tufts.edu:~/project/my_experiment/
\end{lstlisting}

\begin{tipbox}
\textbf{推荐工作流}：每次修改完代码后，用 \texttt{rsync} 同步到集群。因为 rsync 只传改动的文件，所以即使项目很大也很快。
\end{tipbox}

\subsection{场景三：传输很大的数据集（用 Globus）}

如果你需要传输几十 GB 甚至 TB 级别的数据，\texttt{scp} 和 \texttt{rsync} 都不太合适——它们走 SSH 通道，速度有限，而且长时间传输容易断。

Tufts 推荐使用 \textbf{Globus}，一个专门用于科研数据传输的平台：

\begin{enumerate}[nosep]
  \item 访问 \url{https://www.globus.org/}，用 Tufts 账号登录
  \item 设置两个``端点''（Endpoint）：一个是你的电脑（安装 Globus Connect Personal），一个是 Tufts HPC 的 Collection
  \item 在网页界面上选择源和目标文件，点击传输
  \item Globus 在后台完成传输，支持断点续传，\textbf{不需要保持网络连接}
\end{enumerate}

\begin{tipbox}
Globus 的关键优势：``Fire and Forget''——发起传输后你可以关掉浏览器，它会在后台自动完成，完成后邮件通知你。使用 Globus 不需要 VPN。
\end{tipbox}

\subsection{场景四：小文件快速上传（用 OnDemand）}

如果只是传个几 MB 的小文件，用 OnDemand 网页界面最方便：

\begin{enumerate}[nosep]
  \item 打开 \url{https://ondemand-prod.pax.tufts.edu/}
  \item 点击 Files $\rightarrow$ Home Directory
  \item 使用 Upload / Download 按钮
\end{enumerate}

\begin{warnbox}
OnDemand 文件传输限制为\textbf{每个文件不超过 976MB}。大文件请用 rsync 或 Globus。
\end{warnbox}

\subsection{文件传输方式总结}

\begin{table}[h]
\centering
\begin{tabular}{llll}
\toprule
场景 & 推荐工具 & 优势 & 限制 \\
\midrule
几个小文件 & \texttt{scp} & 简单直接 & 不支持续传 \\
项目目录同步 & \texttt{rsync} & 增量、续传 & 需要 SSH 连接 \\
超大数据集 & Globus & 后台传输、不限大小 & 需要配置端点 \\
临时小文件 & OnDemand 网页 & 图形化、零配置 & 单文件 $<$ 976MB \\
\bottomrule
\end{tabular}
\end{table}
       % 文件传输：scp、rsync、Globus
%=============================================================================
\section{软件环境：怎样使用程序和包？}
%=============================================================================

在你自己的电脑上，你可能直接 \texttt{pip install} 就行了。但在 HPC 上，软件管理方式不太一样。

\begin{warnbox}
HPC 集群上你\textbf{没有管理员权限}，不能用 \texttt{apt-get}、\texttt{yum} 等系统包管理器安装软件。所有软件要么通过 Module 系统加载，要么通过 Conda/pip 安装到你自己的目录里。

另外，\textbf{不要在登录节点上跑任何计算或安装操作}。即使是创建 Conda 环境、安装 Python 包，也应该先用 \texttt{srun} 进入一个交互式计算节点再操作。
\end{warnbox}

\subsection{Module 系统：集群的``软件开关''}

HPC 集群上预装了大量软件（Python、CUDA、GCC、R 等等），但它们默认\textbf{不在你的环境变量里}。你需要用 \texttt{module load} 命令来``激活''它们。

为什么要这样设计？因为不同用户可能需要不同版本的同一个软件。Module 系统让每个人可以精确控制自己使用哪个版本，互不干扰。

\textbf{常用 module 命令：}

\begin{lstlisting}[language=bash]
# 查看有哪些软件可用（当前环境）
module av

# 搜索特定软件（搜索全部，包括有依赖关系的）
module spider python
module spider cuda

# 加载软件
module load miniforge/25.3.0  # 加载 Conda（推荐）
module load cuda/12.9.0       # 加载指定版本的 CUDA

# 查看当前加载了什么
module list

# 卸载某个软件
module unload miniforge

# 清空所有已加载的软件
module purge
\end{lstlisting}

\begin{tipbox}
\texttt{module av} 只显示当前可以直接加载的软件。\texttt{module spider} 搜索范围更广，会告诉你某个软件是否存在、需要先加载哪些依赖。找软件时优先用 \texttt{module spider}。
\end{tipbox}

\subsection{Python 环境管理：Conda（官方推荐方式）}

Tufts HPC 官方推荐使用 \textbf{Conda} 来管理 Python 环境和安装包。集群提供了两个 Conda module：

\begin{itemize}[nosep]
  \item \texttt{miniforge/25.3.0}（推荐）— 更轻量
  \item \texttt{anaconda/2025.06.0} — 更完整，预装更多包
\end{itemize}

\subsubsection{第一步：配置存储路径（只需做一次）}

Conda 环境和包缓存会占用大量空间。如果用默认路径，会很快撑满 Home 目录的 30GB 配额。\textbf{必须先把 Conda 的存储路径指向 Research 存储}：

\begin{lstlisting}[language=bash]
# 先创建目录
mkdir -p /cluster/tufts/your_lab/$USER/condaenv
mkdir -p /cluster/tufts/your_lab/$USER/condapkg

# 配置 conda 使用这些目录
conda config --append envs_dirs /cluster/tufts/your_lab/$USER/condaenv/
conda config --append pkgs_dirs /cluster/tufts/your_lab/$USER/condapkg/
\end{lstlisting}

这会在你的 Home 目录下创建一个 \texttt{\textasciitilde/.condarc} 配置文件。之后创建的所有环境和下载的包都会存储在 Research 存储里，而不是占用 Home 目录的配额。

\begin{warnbox}
如果你在 \textbf{Home 目录}下（而不是 Research 存储里）看到 \texttt{condaenv} 和 \texttt{condapkg} 文件夹，说明 Conda 的存储路径配置到了 Home 目录——这会很快撑满 30GB 配额。应该把它们迁移到 Research 存储：

\begin{lstlisting}[language=bash]
# 把内容搬到 Research 存储
mv ~/condaenv/* /cluster/tufts/your_lab/$USER/condaenv/
mv ~/condapkg/* /cluster/tufts/your_lab/$USER/condapkg/
rmdir ~/condaenv ~/condapkg

# 然后更新 .condarc（参考上面的 conda config 命令）
\end{lstlisting}
\end{warnbox}

\subsubsection{第二步：创建和使用环境}

\begin{lstlisting}[language=bash]
# 进入交互式计算节点（不要在登录节点上操作！）
srun -p batch -t 0-2:00:00 -n 2 --mem=4g --pty bash

# 加载 Conda
module load miniforge/25.3.0

# 创建一个新环境
conda create -n rl_env python=3.11 -y

# 激活环境
conda activate rl_env

# 安装包（优先用 conda install，找不到的包再用 pip）
conda install pytorch gymnasium numpy pandas -c conda-forge
# 或者用 pip（conda 里找不到的包）
pip install some-rare-package
\end{lstlisting}

\begin{warnbox}
\textbf{不要混用 \texttt{conda install} 和 \texttt{pip install}}——这可能破坏环境。最佳实践：先用 \texttt{conda install} 装所有能装的包，最后再用 \texttt{pip install} 补充 conda 里找不到的包。
\end{warnbox}

\subsubsection{在 SLURM 脚本中使用 Conda 环境}

\begin{lstlisting}[language=slurm]
#!/bin/bash
#SBATCH -J my_job
#SBATCH --time=00-06:00:00
#SBATCH -p gpu
#SBATCH --gres=gpu:1

module load miniforge/25.3.0
module load cuda/12.9.0
conda activate rl_env

python train.py --lr 0.001
\end{lstlisting}

\subsection{关于 uv：需要额外安装}

如果你在本地习惯用 \texttt{uv} 管理 Python 项目，在 HPC 上也可以用，但需要注意：\textbf{uv 不是集群预装的软件}，没有对应的 module，你需要自己安装。

\begin{lstlisting}[language=bash]
# 进入交互式计算节点
srun -p batch -t 0-2:00:00 -n 2 --mem=4g --pty bash

# 方法：先加载 Conda，用 pip 安装 uv 到用户目录
module load miniforge/25.3.0
pip install --user uv
\end{lstlisting}

安装后 \texttt{uv} 会被放在 \texttt{\textasciitilde/.local/bin/} 下。你需要确保这个路径在 \texttt{PATH} 里：

\begin{lstlisting}[language=bash]
# 检查 uv 是否可用
which uv

# 如果找不到，手动添加 PATH（也可以写入 ~/.bashrc）
export PATH="$HOME/.local/bin:$PATH"
\end{lstlisting}

\begin{warnbox}
以下事项需要注意：
\begin{itemize}[nosep]
  \item \texttt{uv} 不是 Tufts 官方支持的工具，遇到问题无法找 HPC 团队帮忙
  \item 计算节点可能没有外网访问权限，\texttt{uv sync} 需要下载包时可能会失败
  \item 如果外网受限，建议在登录节点完成 \texttt{uv sync}（仅限安装包，不要跑计算），或改用 Conda
\end{itemize}
\end{warnbox}

\subsubsection{正确的使用方式：先手动 sync，脚本里直接用 .venv}

\textbf{关键原则：不要在 SLURM 脚本里写 \texttt{uv sync}。}

原因：如果你提交了一个 Array Job（比如 50 个任务），它们可能同时启动、同时执行 \texttt{uv sync}。每个任务都会尝试下载完整的依赖（比如 PyTorch 就好几个 GB），50 个任务同时下载 = 瞬间吃掉几十 GB 的缓存空间，直接撑爆你的存储配额。

正确做法是\textbf{分两步走}：

\textbf{第一步：手动登录计算节点，运行一次 \texttt{uv sync}}

\begin{lstlisting}[language=bash]
# 开一个交互式节点
srun -p batch -t 0-1:00:00 --mem=4g --pty bash

# 进入项目目录，同步依赖（只需做一次，或依赖变化时重做）
cd ~/project/my_experiment
uv sync
\end{lstlisting}

这会在项目目录下创建 \texttt{.venv/}，里面包含所有依赖。因为所有节点共享同一个文件系统，这个 \texttt{.venv} 在任何计算节点上都能直接使用。

\texttt{uv sync} 会在 \texttt{\textasciitilde/.cache/uv/} 里缓存下载的包，下次 sync 时如果依赖没变就不需要重新下载。但这个缓存可能比较大（PyTorch 等大包会占几十 GB），如果 Home 配额吃紧，可以用 \texttt{rm -rf \textasciitilde/.cache/uv} 清理。

\textbf{第二步：SLURM 脚本里直接调用 .venv 中的 Python}

\begin{lstlisting}[language=slurm]
#!/bin/bash
#SBATCH -J my_job
#SBATCH --time=00-06:00:00
#SBATCH -p gpu
#SBATCH --gres=gpu:1

module load cuda/12.9.0

cd ~/project/my_experiment
# 直接用 .venv 里的 python，不要 uv sync！
.venv/bin/python train.py
\end{lstlisting}

注意最后一行：\texttt{.venv/bin/python} 直接调用虚拟环境里的 Python 解释器，不需要 \texttt{uv run} 或 \texttt{uv sync}，也不需要 \texttt{source .venv/bin/activate}。


\begin{tipbox}
推荐的 uv 工作流（完整版）：
\begin{enumerate}[nosep]
  \item 在本地开发，用 \texttt{uv add} 管理依赖
  \item 用 \texttt{rsync} 把项目（包括 \texttt{pyproject.toml} 和 \texttt{uv.lock}）同步到集群
  \item 在集群上手动 \texttt{srun} 进交互式节点，运行 \texttt{uv sync}（只需一次，或依赖变化时重做）
  \item SLURM 脚本里用 \texttt{.venv/bin/python} 直接运行，\textbf{不要写 \texttt{uv sync}}
\end{enumerate}
如果遇到网络问题导致 \texttt{uv sync} 失败，建议回退到 Conda 方案。
\end{tipbox}
            % 软件环境：Module、Conda、uv

% ---- 第三部分：存储与资源管理（定期关注）----
%=============================================================================
\section{存储系统：你的文件存在哪里？}
%=============================================================================

理解存储结构是使用 HPC 的基础。集群有一个 3PB 的共享存储系统，\textbf{所有登录节点和计算节点都能访问}。

\subsection{登录后看到的文件都是谁的？}

这是很多新手的困惑：登录集群后 \texttt{ls} 一下，看到一堆文件夹，不确定哪些是自己的、哪些是别人的、能不能删。

\textbf{简单的答案是：你登录后看到的所有文件都是你自己的。}

原因如下：
\begin{itemize}[nosep]
  \item 你 SSH 登录后，默认落在你的 \textbf{Home 目录}（\texttt{/cluster/home/your\_utln}）
  \item 这个目录是\textbf{你的私人空间}，只有你能看到和修改（除非你主动改了权限）
  \item 其他用户的 Home 目录在 \texttt{/cluster/home/别人的utln}，和你的完全隔离
  \item 即使你 \texttt{srun} 进入计算节点，看到的文件也一样——因为所有节点共享同一个文件系统，你还是在你自己的 Home 目录里
\end{itemize}

你可以用以下命令确认自己在哪里、文件归谁：

\begin{lstlisting}[language=bash]
pwd                  # 查看当前目录（应该显示 /cluster/home/your_utln）
ls -la               # 查看文件详情，第 3 列是文件所有者
whoami               # 确认当前用户身份
\end{lstlisting}

\texttt{ls -la} 的输出中，第 3 列显示文件所有者，第 4 列显示所属组：

\begin{lstlisting}
drwxr-xr-x  3 your_utln your_utln 4096 Mar 23 10:00 condaenv
drwxr-xr-x  5 your_utln your_utln 4096 Mar 23 10:00 condapkg
                ^^^^^^^^^
                这里显示的是文件所有者，全是你自己
\end{lstlisting}

\begin{tipbox}
\textbf{常见的``遗留文件''及其来源：}
\begin{itemize}[nosep]
  \item \texttt{condaenv} / \texttt{condapkg} — Conda 的环境和包缓存目录。\textbf{如果它们在 Home 目录下，说明 Conda 路径配置不对}——应该迁移到 Research 存储里（见软件环境章节），否则会很快撑满 30GB 配额
  \item \texttt{ondemand} — OnDemand 网页界面自动创建的工作目录，\textbf{不要删}
  \item \texttt{.ipynb} 文件 — 你之前用 Jupyter Notebook 创建的，确认不需要可以删
  \item 软件名目录（如 \texttt{Wolfram Mathematica}）— 你之前使用过该软件时留下的，确认不需要可以删
\end{itemize}
这些都是\textbf{你自己之前操作时产生的}，不是别人的文件，放心处理。
\end{tipbox}

\subsection{两类存储空间}

\begin{table}[h]
\centering
\begin{tabular}{lp{3cm}p{3cm}p{4cm}}
\toprule
类型 & 路径 & 配额 & 用途 \\
\midrule
\textbf{Home 目录} & \texttt{/cluster/home/\linebreak your\_utln} & 30 GB（固定） & 个人配置、小型代码 \\
\textbf{Research 存储} & \texttt{/cluster/tufts/\linebreak lab\_name/your\_utln} & 至少 50 GB（可扩展） & 实验数据、大型项目 \\
\bottomrule
\end{tabular}
\end{table}

Research 存储和 Home 目录的归属逻辑类似：\texttt{/cluster/tufts/lab\_name/your\_utln} 这个子目录是你的个人空间。但 \texttt{/cluster/tufts/lab\_name/} 这一层是整个实验室共享的，你可能能看到同事的目录（取决于权限设置）。

\begin{warnbox}
Home 目录只有 30GB，很容易被 conda 环境和数据集撑满。\textbf{建议把实验项目、数据集、conda 环境都放在 Research 存储里。}

查看用量：\texttt{df -H /cluster/tufts/your\_lab\_name}
\end{warnbox}

\begin{tipbox}
如果你想确保自己的 Home 目录不被其他用户查看，可以运行：
\begin{lstlisting}[language=bash]
chmod 700 ~
\end{lstlisting}
这会把权限设为``只有自己可读写执行''。
\end{tipbox}

\subsection{集群禁止存放受限数据}

根据 Tufts HPC 政策，集群上\textbf{不允许存放受限数据}（Restricted Data），包括 HIPAA、FERPA 等受保护的信息。更多使用规则见附录 D。
             % 存储：Home vs Research、配额、文件归属
%=============================================================================
\section{目录管理：怎样组织你的文件？}
%=============================================================================

集群不像你的个人电脑——它的存储有配额限制，而且你可能同时跑几十个实验。花几分钟规划好目录结构，会在后续省掉很多麻烦。

\subsection{推荐的目录布局}

\begin{lstlisting}
/cluster/home/your_utln/              # Home 目录（30GB）
  .condarc                            # Conda 配置（指向 Research 存储）
  .bashrc                             # Shell 配置
  privatemodules/                     # 自定义 module（如果需要的话）

/cluster/tufts/your_lab/your_utln/    # Research 存储（50GB+）
  condaenv/                           # Conda 环境（按官方推荐配置到这里）
  condapkg/                           # Conda 包缓存
  projects/                           # 所有实验项目
    rl_matrix/                        # 项目 1
      train.py
      configs/
      scripts/
      results/
      logs/
    nlp_finetune/                     # 项目 2
      ...
  datasets/                           # 共享数据集（避免每个项目复制一份）
    mujoco_data/
    imagenet_subset/
\end{lstlisting}

\subsection{核心原则}

\begin{description}[style=nextline,leftmargin=2em]
  \item[Home 目录只放配置文件]
    \texttt{\textasciitilde/.condarc}、\texttt{\textasciitilde/.bashrc} 等配置文件很小，放 Home 没问题。但代码、数据、Conda 环境都应该放到 Research 存储里——Home 只有 30GB，一个 PyTorch 环境就能吃掉一半。

  \item[每个项目一个独立目录]
    和你在自己电脑上一样，每个实验项目放在独立的目录里，内部包含代码、配置、脚本、结果、日志。这样项目之间不会互相干扰，也方便用 \texttt{rsync} 同步。

  \item[结果和日志与代码分开]
    在项目目录里用 \texttt{results/} 和 \texttt{logs/} 子目录分别存放程序输出和 SLURM 日志。这样清理日志时不会误删代码，下载结果时也不用把整个项目都拉下来。

  \item[数据集单独存放、项目间共享]
    如果多个项目用同一份数据集，不要复制多份。在 Research 存储里建一个 \texttt{datasets/} 目录，各项目通过路径引用它。

  \item[定期清理]
    实验跑完、论文发了之后，把不需要的 checkpoint、日志和中间文件删掉释放空间。特别是 SLURM 的 \texttt{.out} 和 \texttt{.err} 日志——一次 Array Job 就能产生上百个日志文件。
\end{description}

\subsection{实用命令}

\begin{lstlisting}[language=bash]
# 查看 Home 目录用了多少空间
du -sh ~

# 查看 Research 存储的配额和用量
df -H /cluster/tufts/your_lab

# 找出占空间最多的目录（排序显示前 10 个）
du -h --max-depth=2 ~ | sort -rh | head -10

# 批量删除旧日志（比如 30 天前的 .out 和 .err 文件）
find logs/ -name "*.out" -mtime +30 -delete
find logs/ -name "*.err" -mtime +30 -delete
\end{lstlisting}

\begin{warnbox}
集群的存储系统\textbf{不提供备份}，但会保留最近 90 天的每日快照（Snapshots），可以恢复最近误删或误改的文件。详情可参考 Tufts 官方的 File Recovery 教程。如果你有重要的实验结果，建议自己也备份一份到本地或云端。
\end{warnbox}
% 目录管理：推荐布局、清理
%=============================================================================
\section{集群资源概览}
%=============================================================================

\subsection{分区（Partition）}

SLURM 把计算节点分成几组，叫做``分区''。不同分区有不同的用途和限制：

\begin{table}[h]
\centering
\begin{tabular}{llp{7cm}}
\toprule
分区 & 时间上限 & 说明 \\
\midrule
\texttt{batch} & 2 天 & 默认分区，仅 CPU。不允许请求 GPU。 \\
\texttt{gpu} & 2 天 & GPU 分区。必须用 \texttt{--gres} 请求 GPU，不允许只跑 CPU 任务。 \\
\texttt{preempt} & 2 天 & 资源最多（包含教授贡献的节点），但作业可能被节点所有者\textbf{抢占}——你的作业会被直接 kill。 \\
\bottomrule
\end{tabular}
\end{table}

\subsection{每用户资源限制}

\begin{table}[h]
\centering
\begin{tabular}{lccc}
\toprule
分区组合 & CPU 核数上限 & 内存上限 & GPU 上限 \\
\midrule
\texttt{batch} + \texttt{gpu} & 250 & 5,000 GB & 10 \\
\texttt{preempt} & 1,000 & 8,000 GB & 20 \\
\bottomrule
\end{tabular}
\end{table}

\begin{tipbox}
可以同时指定多个分区做 fallback：\texttt{\#SBATCH -p gpu,preempt}。这样如果 \texttt{gpu} 分区没资源，会自动尝试 \texttt{preempt}。
\end{tipbox}

\subsection{可用 GPU}

集群有多种型号的 GPU：

\begin{table}[h]
\centering
\begin{tabular}{lcl}
\toprule
GPU 型号 & 显存 & 适用场景 \\
\midrule
T4 & 16 GB & 小模型推理、轻量训练 \\
L40 / L40s & 48 GB & 中等规模训练 \\
A100-40G & 40 GB & 大规模训练 \\
A100-80G & 80 GB & 大模型训练 \\
H200 & 141 GB & 最大规模任务 \\
\bottomrule
\end{tabular}
\end{table}

请求特定型号的 GPU：

\begin{lstlisting}[language=slurm]
#SBATCH --gres=gpu:a100:1       # 请求 1 块 A100（不区分 40G/80G）
#SBATCH --constraint="a100-80G" # 加上这行可以指定 80G 版本
\end{lstlisting}

\begin{warnbox}
除非你的框架支持多 GPU 并行，否则\textbf{不要请求多块 GPU}——你只会浪费资源、延长排队时间。
\end{warnbox}

\subsection{怎样查看集群当前的资源占用情况？}

在提交作业之前，你可能想知道：现在哪个分区有空闲资源？GPU 还剩多少？这可以帮你决定选哪个分区、请求什么类型的 GPU，避免排长队。

\subsubsection{方法一：hpctools（推荐，最直观）}

Tufts 提供了一个叫 \texttt{hpctools} 的辅助工具，用交互式菜单帮你查看各种信息：

\begin{lstlisting}[language=bash]
module load hpctools
hpctools
\end{lstlisting}

运行后会显示一个菜单，输入编号选择功能：

\begin{table}[h]
\centering
\begin{tabular}{cl}
\toprule
选项 & 功能 \\
\midrule
1 & 查看各分区可用的 CPU 核数和内存 \\
2 & 查看各分区可用的 GPU 资源 \\
3 & 查看某日期之后你的已完成作业历史 \\
4 & 查看你当前运行的作业详情 \\
5 & 查看 Research 存储的配额和用量 \\
6 & 查看指定目录的详细空间占用 \\
7 & 查看 Home 目录的空间占用明细 \\
\bottomrule
\end{tabular}
\end{table}

比如你想在提交 GPU 作业前看看还有多少 GPU 可用，选 \texttt{2} 就能看到各分区的 GPU 空闲情况。想检查自己的存储快满了没有，选 \texttt{5} 或 \texttt{7}。

\subsubsection{方法二：sinfo（命令行快速查看）}

如果你不想进交互式菜单，\texttt{sinfo} 可以一行命令看到集群状态：

\begin{lstlisting}[language=bash]
# 查看所有分区的节点状态概览
sinfo

# 只看 gpu 分区，显示空闲/占用/总计
sinfo -p gpu -o "%P %a %D %C"
\end{lstlisting}

输出中 \texttt{\%C} 格式为 \texttt{已分配/空闲/其他/总计} 的 CPU 核数。

查看 GPU 分区中各节点的具体 GPU 空闲情况：

\begin{lstlisting}[language=bash]
# 显示每个节点的 GPU 分配情况
sinfo -p gpu -O NodeList,Gres,GresUsed,CPUsState,FreeMem
\end{lstlisting}

\subsubsection{方法三：OnDemand 网页界面}

如果你偏好图形化界面，OnDemand 提供了两个有用的页面：

\begin{itemize}[nosep]
  \item \textbf{Clusters $\rightarrow$ System Status}：查看整个集群各分区的节点状态
  \item \textbf{Jobs $\rightarrow$ Active Jobs}：查看你自己的作业运行情况，点击单个作业可以看详情，已完成的作业还能在 Resources and I/O 标签下看到 CPU 和内存的效率
\end{itemize}

\subsubsection{查看正在运行的作业的实时资源使用}

如果你的作业已经在跑了，想看它实际占用了多少 CPU 和 GPU：

\begin{lstlisting}[language=bash]
# SSH 到作业所在的节点（先用 squeue --me 找到节点名）
ssh pax006

# 查看你的进程的 CPU 和内存使用
top -u $USER

# 查看 GPU 使用情况（在有 GPU 的节点上）
nvidia-smi
\end{lstlisting}

\texttt{nvidia-smi} 会显示每块 GPU 的算力利用率、显存占用、以及正在使用 GPU 的进程。如果你发现 GPU 利用率很低（比如低于 30\%），可能说明你的程序瓶颈在数据加载而不是计算——可以尝试增加 \texttt{--cpus-per-task} 来加速数据预处理。
           % 集群资源：分区、GPU、监控工具

% ---- 第四部分：背景知识与进阶（读一次，需要时回来查）----
%=============================================================================
\section{SLURM 的核心哲学：理解 HPC 的运行方式}
%=============================================================================

在开始任何操作之前，你需要先理解 HPC 集群和你的笔记本电脑有什么本质区别。

\subsection{你的电脑 vs HPC 集群}

在你自己的电脑上，运行一个 Python 程序很简单：打开终端，输入 \texttt{python train.py}，程序就在你的电脑上跑了。你是这台电脑的唯一用户，所有资源（CPU、内存、GPU）都归你一个人用。

HPC 集群完全不一样。它是一台\textbf{由上百台服务器组成的``超级计算机''}，同时有几十甚至上百个用户在使用。如果大家都直接运行程序，就像几百个人同时在一个厨房里做饭——必然乱套。

所以集群需要一个\textbf{调度系统}来维持秩序：你不能直接跑程序，而是要先``申请资源''，说明你需要多少 CPU、多少内存、用多长时间，然后系统帮你安排。

\subsection{SLURM 是什么？}

\textbf{SLURM}（Simple Linux Utility for Resource Management）就是 Tufts 集群使用的调度系统。你可以把它想象成一个\textbf{餐厅的排号系统}：

\begin{enumerate}[nosep]
  \item \textbf{你告诉服务员}（写一个作业脚本）：``我需要一个 4 人桌、靠窗的位置、用 2 小时''
  \item \textbf{系统排号}（SLURM 把你的请求放进队列）
  \item \textbf{有位置了就安排你入座}（当集群有足够空闲资源时，SLURM 把你分配到某台服务器上）
  \item \textbf{你吃完走人}（程序运行完毕，结果写入文件，资源释放）
\end{enumerate}

整个过程是\textbf{异步的}——你提交作业后可以关掉终端去做别的事，程序会在后台自动运行。

\subsection{两种关键的节点：登录节点 vs 计算节点}

\begin{conceptbox}
当你 SSH 连接到集群时，你降落在\textbf{登录节点}（Login Node）上。登录节点就像一个大厅，只用来做轻量操作：编辑文件、提交作业、查看结果。

真正跑计算的是\textbf{计算节点}（Compute Node）——你通过 SLURM 申请之后才能使用它们。

\textbf{永远不要在登录节点上直接跑训练程序}，这会影响所有其他用户。
\end{conceptbox}

那么具体什么\textbf{可以}在登录节点上做，什么\textbf{不可以}？一个简单的判断标准：

\begin{table}[h]
\centering
\begin{tabular}{lp{8cm}}
\toprule
可以在登录节点做 & 不可以在登录节点做 \\
\midrule
编辑文件（vim, nano 等） & 运行训练脚本（python train.py） \\
\texttt{module load}（只是修改环境变量） & 处理大型数据集 \\
\texttt{sbatch} 提交作业 & 运行消耗大量 CPU/GPU 的程序 \\
\texttt{squeue}、\texttt{sinfo} 查看状态 & 编译大型软件 \\
\texttt{hpctools} 查看资源 & \\
\texttt{rsync}/\texttt{scp} 传文件 & \\
\texttt{conda create} 创建环境* & \\
\bottomrule
\end{tabular}
\end{table}

\noindent *注：创建 Conda 环境和安装包涉及一定计算，官方建议在交互式计算节点上操作，但如果只是小规模安装（几个包），在登录节点上做通常也不会有问题。

Tufts 集群有 3 个登录节点（\texttt{login-p01}, \texttt{login-p02}, \texttt{login-p03}），通过统一入口 \texttt{login-prod.pax.tufts.edu} 自动分配。

\subsection{你的程序到底是怎么``跑起来''的？}

这是很多新手最困惑的地方，让我们一步步拆解：

\begin{enumerate}
  \item \textbf{你的代码和数据}存放在集群的共享存储系统上（一个 3PB 的大硬盘，所有节点都能访问）
  \item 你写一个\textbf{作业脚本}（就是一个 \texttt{.sh} 文件），里面说明：需要什么资源、要运行什么命令
  \item 你用 \texttt{sbatch} 命令提交这个脚本
  \item SLURM 找到一台空闲的计算节点，在上面\textbf{以你的身份}执行这个脚本
  \item 因为所有节点共享同一个存储系统，计算节点\textbf{能直接读取你的代码和数据}——不需要额外复制
  \item 程序的标准输出（\texttt{print} 的内容）被重定向到一个日志文件里
  \item 程序产生的输出文件（模型、结果等）直接写在共享存储上，跑完后你在登录节点就能看到
\end{enumerate}

\begin{tipbox}
关键理解：\textbf{所有节点看到的是同一套文件系统。}你在登录节点上传的文件，计算节点立刻就能读取；计算节点写出的结果，你在登录节点立刻就能下载。不存在``从计算节点复制结果到登录节点''的步骤。
\end{tipbox}
    % SLURM 哲学：为什么需要调度器
%=============================================================================
\section{实战：强化学习实验矩阵}
%=============================================================================

现在我们把前面学到的知识串起来，实现一个完整的实验工作流。

我们的实验是一个三维矩阵：

\begin{itemize}[nosep]
  \item \textbf{维度 1}：多个 RL 环境（CartPole, LunarLander, HalfCheetah, \ldots）
  \item \textbf{维度 2}：多个算法/架构（PPO, SAC, DQN, \ldots）
  \item \textbf{维度 3}：每组跑 $N$ 个随机种子（确保统计显著性）
\end{itemize}

目标：一键提交所有实验，自动记录结果，方便后续分析。

\subsection{项目结构}

先来看整个项目长什么样：

\begin{lstlisting}
project/
  train.py              # 训练脚本，接受命令行参数
  configs/
    envs.txt            # 环境列表，每行一个
    archs.txt           # 架构列表，每行一个
  scripts/
    run_matrix.sh       # SLURM Array Job 脚本
    submit_all.sh       # 提交脚本
  results/              # 程序输出的结果
  logs/                 # SLURM 日志文件
\end{lstlisting}

\texttt{results/} 存放你程序写出的模型和指标，\texttt{logs/} 存放 SLURM 的标准输出和错误日志。两者分开管理比较清晰。

\subsection{配置文件}

用纯文本文件列出实验参数，方便修改，也方便脚本读取：

\begin{lstlisting}[title={\texttt{configs/envs.txt}}]
CartPole-v1
LunarLander-v2
HalfCheetah-v4
Hopper-v4
Walker2d-v4
\end{lstlisting}

\begin{lstlisting}[title={\texttt{configs/archs.txt}}]
PPO
SAC
DQN
TD3
A2C
\end{lstlisting}

\subsection{训练脚本接口}

\texttt{train.py} 需要接受命令行参数，这样 SLURM 脚本才能通过参数控制每次运行的内容：

\begin{lstlisting}[language=python,title={\texttt{train.py}（参数部分）}]
import argparse

parser = argparse.ArgumentParser()
parser.add_argument("--env", type=str, required=True)
parser.add_argument("--arch", type=str, required=True)
parser.add_argument("--seed", type=int, required=True)
parser.add_argument("--output-dir", type=str, default="results")
args = parser.parse_args()

# 结果保存到: results/{env}/{arch}/seed_{seed}/
\end{lstlisting}

关键设计：把环境、算法、种子都做成参数，这样同一个脚本就能覆盖所有实验组合。

%=============================================================================
\subsection{方案一：SLURM Array Job（推荐）}
%=============================================================================

SLURM Array Job 允许你\textbf{一次提交一批结构相同但参数不同的作业}。核心思路很简单：

\begin{enumerate}[nosep]
  \item 把三维矩阵（环境 $\times$ 架构 $\times$ 种子）展开为一个一维的数字索引 \texttt{0, 1, 2, ..., N-1}
  \item SLURM 为每个索引创建一个独立的 Task，通过环境变量 \texttt{SLURM\_ARRAY\_TASK\_ID} 告诉脚本``你是第几号''
  \item 脚本根据这个编号反推出对应的（环境, 架构, 种子）组合
\end{enumerate}

假设有 $E$ 个环境、$A$ 个架构、每组 $N$ 个种子，总任务数 $= E \times A \times N$。

\begin{lstlisting}[language=slurm,title={\texttt{scripts/run\_matrix.sh}}]
#!/bin/bash
#SBATCH -J rl_matrix
#SBATCH --time=00-12:00:00
#SBATCH -p gpu,preempt
#SBATCH -N 1
#SBATCH -n 1
#SBATCH --cpus-per-task=4
#SBATCH --mem=16g
#SBATCH --gres=gpu:1
#SBATCH --output=logs/%x_%A_%a.out
#SBATCH --error=logs/%x_%A_%a.err
#SBATCH --mail-type=END,FAIL
\end{lstlisting}

注意 \texttt{--output} 里用了 \texttt{\%A}（主 Job ID）和 \texttt{\%a}（Task 序号），这样每个 Task 有独立的日志文件。

\begin{lstlisting}[language=slurm]
# ---- 配置 ----
NUM_SEEDS=5

# 读取环境和架构列表到数组
mapfile -t ENVS < configs/envs.txt
mapfile -t ARCHS < configs/archs.txt

NUM_ENVS=${#ENVS[@]}
NUM_ARCHS=${#ARCHS[@]}
\end{lstlisting}

\texttt{mapfile} 把文本文件逐行读入 Bash 数组。这样我们可以用索引来访问每个元素。

\begin{lstlisting}[language=slurm]
# ---- 从 SLURM_ARRAY_TASK_ID 解码实验参数 ----
# 把一维索引还原为三维坐标，类似于把数组下标转换为行列号
TASK_ID=${SLURM_ARRAY_TASK_ID}

SEED=$((TASK_ID % NUM_SEEDS))          # 最内层：种子
TASK_ID=$((TASK_ID / NUM_SEEDS))
ARCH_IDX=$((TASK_ID % NUM_ARCHS))      # 中间层：架构
ENV_IDX=$((TASK_ID / NUM_ARCHS))       # 最外层：环境

ENV=${ENVS[$ENV_IDX]}
ARCH=${ARCHS[$ARCH_IDX]}

echo "=== Task ${SLURM_ARRAY_TASK_ID}: env=${ENV}, arch=${ARCH}, seed=${SEED} ==="
\end{lstlisting}

这段代码的逻辑和``把一个一维数组下标转成多维下标''完全一样（取余、整除），如果你学过数组的线性化存储就会觉得很熟悉。

\begin{lstlisting}[language=slurm]
# ---- 加载环境并运行 ----
module load cuda/12.9.0
module load miniforge/25.3.0
conda activate rl_env

python train.py \
  --env "$ENV" \
  --arch "$ARCH" \
  --seed "$SEED" \
  --output-dir "results/${ENV}/${ARCH}/seed_${SEED}"
\end{lstlisting}

\subsubsection{提交脚本}

提交脚本负责计算总任务数，然后一条命令提交所有任务：

\begin{lstlisting}[language=bash,title={\texttt{scripts/submit\_all.sh}}]
#!/bin/bash
# 计算总任务数
NUM_ENVS=$(wc -l < configs/envs.txt)
NUM_ARCHS=$(wc -l < configs/archs.txt)
NUM_SEEDS=5

TOTAL=$((NUM_ENVS * NUM_ARCHS * NUM_SEEDS))
MAX_INDEX=$((TOTAL - 1))

echo "Submitting ${TOTAL} tasks (${NUM_ENVS} envs x ${NUM_ARCHS} archs x ${NUM_SEEDS} seeds)"

# 创建日志目录（必须提前存在，否则 SLURM 写日志会失败）
mkdir -p logs results

# 提交 Array Job
sbatch --array=0-${MAX_INDEX} scripts/run_matrix.sh
\end{lstlisting}

以 5 环境 $\times$ 5 架构 $\times$ 5 种子为例，共 125 个任务，一条命令全部提交。

\subsubsection{控制并发数}

集群对每用户有 GPU 数量限制。用 \texttt{\%} 符号限制同时运行的 Array Task 数：

\begin{lstlisting}[language=bash]
# 最多同时跑 10 个任务（因为 gpu 分区每用户限 10 块 GPU）
sbatch --array=0-124%10 scripts/run_matrix.sh
\end{lstlisting}

如果你同时用了 \texttt{preempt} 分区，上限可以到 20。

%=============================================================================
\subsection{方案二：嵌套循环提交}
%=============================================================================

如果不同实验需要不同的资源配置（比如 MuJoCo 环境需要更多内存），可以用脚本循环提交独立作业，对每组实验单独设置参数：

\begin{lstlisting}[language=bash,title={\texttt{scripts/submit\_loop.sh}}]
#!/bin/bash
NUM_SEEDS=5

while IFS= read -r ENV; do
  while IFS= read -r ARCH; do
    for SEED in $(seq 0 $((NUM_SEEDS - 1))); do
      # 根据环境选择资源配置
      if [[ "$ENV" == *"Cheetah"* || "$ENV" == *"Hopper"* || "$ENV" == *"Walker"* ]]; then
        MEM="32g"
        TIME="01-00:00:00"
      else
        MEM="8g"
        TIME="00-04:00:00"
      fi

      sbatch \
        -J "rl_${ENV}_${ARCH}_s${SEED}" \
        --time=${TIME} \
        -p gpu,preempt \
        --gres=gpu:1 \
        --cpus-per-task=4 \
        --mem=${MEM} \
        --output="logs/${ENV}_${ARCH}_s${SEED}.out" \
        --error="logs/${ENV}_${ARCH}_s${SEED}.err" \
        --wrap="module load cuda/12.9.0 && module load miniforge/25.3.0 && \
                conda activate rl_env && \
                python train.py --env ${ENV} --arch ${ARCH} --seed ${SEED} \
                --output-dir results/${ENV}/${ARCH}/seed_${SEED}"
    done
  done < configs/archs.txt
done < configs/envs.txt
\end{lstlisting}

\texttt{sbatch --wrap} 可以直接传入命令字符串作为作业内容，不需要单独写作业脚本文件。适合逻辑简单的场景。
       % 实战：RL 实验矩阵（Array Job）
%=============================================================================
\section{实用技巧}
%=============================================================================

\subsection{利用 preempt 分区跑更多实验}

\texttt{preempt} 分区的资源上限更高（20 GPU vs 10 GPU），但作业可能被节点所有者的任务\textbf{抢占}——SLURM 会在约 30 秒内 kill 你的进程。

应对策略是\textbf{让你的程序支持断点续训}：

\begin{lstlisting}[language=python,title={在 train.py 中添加 checkpoint 逻辑}]
import os, torch

ckpt_path = f"{args.output_dir}/checkpoint.pt"

# 启动时检查是否有已保存的进度
if os.path.exists(ckpt_path):
    checkpoint = torch.load(ckpt_path)
    model.load_state_dict(checkpoint["model"])
    start_epoch = checkpoint["epoch"]
    print(f"Resumed from epoch {start_epoch}")
else:
    start_epoch = 0

# 训练过程中定期保存
for epoch in range(start_epoch, total_epochs):
    train_one_epoch(model, ...)
    if epoch % 10 == 0:  # 每 10 个 epoch 保存一次
        torch.save({"model": model.state_dict(), "epoch": epoch}, ckpt_path)
\end{lstlisting}

这样即使作业被抢占，重新提交后也能从上次的进度继续，而不是从头开始。

\subsection{用邮件通知监控实验}

\begin{lstlisting}[language=slurm]
#SBATCH --mail-type=END,FAIL
#SBATCH --mail-user=your.email@tufts.edu
\end{lstlisting}

对于 Array Job（一次提交上百个任务），建议只在失败时通知，否则你的邮箱会被淹没：

\begin{lstlisting}[language=slurm]
#SBATCH --mail-type=FAIL
\end{lstlisting}

\subsection{汇总实验结果}

所有实验跑完后，用一个简单的 Python 脚本汇总各组的结果：

\begin{lstlisting}[language=python,title={\texttt{collect\_results.py}}]
import os, json, pandas as pd

rows = []
for env in open("configs/envs.txt"):
    env = env.strip()
    for arch in open("configs/archs.txt"):
        arch = arch.strip()
        for seed in range(5):
            result_file = f"results/{env}/{arch}/seed_{seed}/metrics.json"
            if os.path.exists(result_file):
                with open(result_file) as f:
                    metrics = json.load(f)
                rows.append({"env": env, "arch": arch, "seed": seed,
                             **metrics})

df = pd.DataFrame(rows)
summary = df.groupby(["env", "arch"]).agg(
    reward_mean=("final_reward", "mean"),
    reward_std=("final_reward", "std"),
).reset_index()
summary.to_csv("results/summary.csv", index=False)
print(summary.to_string())
\end{lstlisting}

\subsection{检查哪些实验失败了}

\begin{lstlisting}[language=bash]
# 查看有错误输出的日志
grep -l "Error\|Traceback\|CANCELLED" logs/*.err

# 查看哪些结果文件缺失
python -c "
import os
envs = open('configs/envs.txt').read().split()
archs = open('configs/archs.txt').read().split()
for e in envs:
    for a in archs:
        for s in range(5):
            p = f'results/{e}/{a}/seed_{s}/metrics.json'
            if not os.path.exists(p):
                print(f'MISSING: {e} / {a} / seed {s}')
"
\end{lstlisting}
                % 实用技巧：preempt、断点续训、邮件通知
%=============================================================================
\section{完整工作流总结}
%=============================================================================

\begin{enumerate}
  \item \textbf{本地准备}：开发代码，准备好 \texttt{train.py}、配置文件和 SLURM 脚本

  \item \textbf{上传到集群}：
\begin{lstlisting}[language=bash]
rsync -azP --exclude='.venv' --exclude='__pycache__' \
  ./my_project/ your_utln@login-prod.pax.tufts.edu:~/project/
\end{lstlisting}

  \item \textbf{SSH 登录，设置环境}：
\begin{lstlisting}[language=bash]
ssh your_utln@login-prod.pax.tufts.edu
module load miniforge/25.3.0 && module load cuda/12.9.0
conda activate rl_env
\end{lstlisting}

  \item \textbf{交互式调试}——先跑一个任务确认一切正常：
\begin{lstlisting}[language=bash]
srun -p gpu -t 0-1:00:00 --gres=gpu:1 --mem=8g --cpus-per-task=4 --pty bash
python train.py --env CartPole-v1 --arch PPO --seed 0 --output-dir test_run
exit
\end{lstlisting}

  \item \textbf{用 seff 检查资源使用}，调整请求量

  \item \textbf{批量提交所有实验}：
\begin{lstlisting}[language=bash]
bash scripts/submit_all.sh
\end{lstlisting}

  \item \textbf{监控进度}：
\begin{lstlisting}[language=bash]
squeue --me                           # 查看状态
watch -n 60 'squeue --me | wc -l'     # 持续监控剩余任务数
\end{lstlisting}

  \item \textbf{下载结果}：
\begin{lstlisting}[language=bash]
# 在本地运行
rsync -azP your_utln@login-prod.pax.tufts.edu:~/project/results/ ./results/
python collect_results.py
\end{lstlisting}
\end{enumerate}
    % 完整工作流总结
%=============================================================================
\section{附录 A：常用 SLURM 命令速查}
%=============================================================================

\begin{table}[h]
\centering
\begin{tabular}{ll}
\toprule
命令 & 功能 \\
\midrule
\texttt{sbatch script.sh} & 提交批处理作业 \\
\texttt{squeue --me} & 查看自己的作业 \\
\texttt{scancel JOBID} & 取消指定作业 \\
\texttt{scancel -u \$USER} & 取消自己的所有作业 \\
\texttt{seff JOBID} & 查看已完成作业的 CPU/内存效率 \\
\texttt{sacct -j JOBID} & 查看作业详细信息 \\
\texttt{scontrol show job JOBID} & 查看作业完整配置 \\
\texttt{sinfo} & 查看集群节点状态 \\
\texttt{nvidia-smi} & 查看 GPU 使用情况（在计算节点上） \\
\bottomrule
\end{tabular}
\end{table}
%=============================================================================
\section{附录 B：常用 \texttt{\#SBATCH} 指令速查}
%=============================================================================

\begin{table}[h]
\centering
\small
\begin{tabular}{lp{9cm}}
\toprule
指令 & 说明 \\
\midrule
\texttt{\#SBATCH -J name} & 作业名称 \\
\texttt{\#SBATCH --time=DD-HH:MM:SS} & 最大运行时间（所有公共分区上限 2 天） \\
\texttt{\#SBATCH -p partition} & 分区（可逗号分隔多个，如 \texttt{gpu,preempt}） \\
\texttt{\#SBATCH -N 1} & 节点数（通常设 1） \\
\texttt{\#SBATCH -n 1} & 任务数 \\
\texttt{\#SBATCH --cpus-per-task=4} & 每任务 CPU 核数 \\
\texttt{\#SBATCH --mem=16g} & 总内存 \\
\texttt{\#SBATCH --gres=gpu:type:N} & GPU 请求（类型和数量） \\
\texttt{\#SBATCH --constraint="..."} & 硬件约束（如 \texttt{a100-80G}） \\
\texttt{\#SBATCH --array=0-99\%10} & Array Job，\%后为最大并发数 \\
\texttt{\#SBATCH --output=\%x.\%j.out} & 标准输出文件 \\
\texttt{\#SBATCH --error=\%x.\%j.err} & 标准错误文件 \\
\texttt{\#SBATCH --mail-type=FAIL} & 邮件通知类型 \\
\texttt{\#SBATCH --mail-user=email} & 通知邮箱 \\
\bottomrule
\end{tabular}
\end{table}
%=============================================================================
\section{附录 C：致谢声明}
%=============================================================================

如果你的研究使用了 Tufts HPC 集群，请在论文中加入致谢：

\begin{quote}
\textit{The authors acknowledge the Tufts University High Performance Compute Cluster (\url{https://it.tufts.edu/high-performance-computing}) which was utilized for the research reported in this paper.}
\end{quote}

%=============================================================================
\section{附录 D：使用规则——初学者容易踩的坑}
\label{sec:policy}
%=============================================================================

HPC 集群受 Tufts \textit{Use of Information Systems Policy} 约束（完整版：\url{https://go.tufts.edu/hpc_aup}）。下面不列那些显而易见的规则，只说\textbf{初学者容易犯的错误}：

\begin{description}[style=nextline,leftmargin=2em]

  \item[不要把账号密码分享给同事]
    即使是同一个实验室的人，也不能共用账号。每个人必须用自己的 UTLN 登录。如果你的同事需要访问集群，让他们自己申请账号。

  \item[不要在集群上存放受限数据]
    HIPAA（医疗数据）、FERPA（学生记录）等受保护的信息\textbf{禁止}存放在 HPC 集群上。如果你不确定你的数据是否属于受限类别，先联系 \texttt{tts-research@tufts.edu} 确认。

  \item[不要在登录节点上跑计算]
    在登录节点上直接跑 \texttt{python train.py} 会影响所有其他用户，管理员可能会直接终止你的进程。

  \item[不要请求远超实际需要的资源]
    ``反正是免费的，我多请求一点''——这个想法是错的。你请求的资源在作业运行期间是\textbf{独占的}，别人无法使用。请求 8 块 GPU 但只用了 1 块，意味着 7 块 GPU 白白浪费。用 \texttt{seff} 检查实际用量，下次按需请求。

  \item[注意软件许可证]
    集群上有些软件\textbf{不是开源的}（如 MATLAB、Mathematica），使用时需要遵守对应的许可协议。违反许可协议算违反使用政策。

  \item[不要用 HPC 做商业用途]
    集群只能用于 Tufts 相关的学术和研究活动。不能用来给外部公司跑计算、挖矿、或者任何商业盈利活动。

  \item[毕业或离职前备份你的数据]
    你离开 Tufts 后，账号访问权限会被取消。\textbf{提前}把需要的数据下载到本地或其他存储，不要等到最后一天。

  \item[你在集群上的活动不是完全私密的]
    Tufts 不会日常监控你的使用，但系统会记录日志。如果有合理理由怀疑违规，管理员可以\textbf{不通知你的情况下}查看你的活动记录和文件。

\end{description}

\begin{tipbox}
以上规则的核心精神其实就一条：\textbf{这是一个共享资源，你的行为会影响其他人。}按需使用、不浪费、不越界，就不会有问题。
\end{tipbox}
            % 附录：命令速查、SBATCH 指令、致谢

\end{document}
